\documentclass[11pt]{article}
\usepackage[a4paper,margin=26mm]{geometry}
\usepackage{amsmath,amssymb}
\usepackage[hidelinks]{hyperref}
\setlength{\emergencystretch}{3em}
\title{Conjecture 00000000721: negative factorial powers do not give an uncountable $p$-adic family}
\author{lizaixi01, with OpenAI Codex assistance}
\date{22 October 2026}
\begin{document}
\maketitle
\section{The stated series and its ambient field}
The source claims that the set
\[
 \left\{\sum_{n\geq1}a_n p^{-n!}:a_n\in\{0,\ldots,p-1\}\right\}
\]
contains an algebraically independent subfamily of cardinality
$2^{\aleph_0}$ in $\overline{\mathbf Q}_p$.
We retain the negative exponent exactly as stated and use the usual
extension of the $p$-adic absolute value, normalized by $|p|_p=p^{-1}$.
At $p=2$, every convergent series in this set is a finite sum and is
rational. Its value set is therefore countable, which excludes the
required cardinality regardless of algebraic independence.

The ambient field in the proof is the actual algebraic closure
$K=\overline{\mathbf Q}_2$ of $\mathbf Q_2$, equipped with its extended
$2$-adic norm. Completeness is unnecessary: we consider series that
actually converge in $K$. Changing the sign in the exponent would
produce a different claim, to which this argument does not apply.

\section{Convergence forces finite support}
Let $a_n\in\{0,1\}$. If $a_n\ne0$, then $a_n=1$ and
\[
 |a_n2^{-n!}|_2=2^{n!}\geq1.                         \tag{1}
\]
In every normed additive group, convergence of the sequence of partial
sums implies that the individual terms tend to zero: subtract the
partial sum ending at $n-1$ from the one ending at $n$.
It follows that eventually $|a_n2^{-n!}|_2<1$.
By (1), all sufficiently late digits are zero.

Consequently there is $N$ such that every convergent value is
\[
 \sum_{n=1}^{N} a_n 2^{-n!}\in\mathbf Q.
\]
The rational numbers are countable, and the field embedding
$\mathbf Q\hookrightarrow K$ carries a countable set to a countable
set. Thus the full set of actual series values is countable. Every
subset is countable, whereas $2^{\aleph_0}>\aleph_0$ by Cantor's theorem.
The specified continuum-sized subfamily cannot exist.

\section{The two meanings of series convergence are covered}
The code first proves the finite-support statement for Lean's
\texttt{HasSum} predicate. It then explicitly treats the ordinary
sequence of forward partial sums
\[
 s_N=\sum_{j=0}^{N-1} a_{j+1}2^{-(j+1)!}.
\]
If $s_N\to x$, then $s_{N+1}-s_N\to0$. Applying (1) to the terms shows
that the digits eventually vanish. The resulting finite sum has a
\texttt{HasSum} proof, and uniqueness of the partial-sum limit identifies
its value with $x$. Conversely, \texttt{HasSum} implies convergence of
these forward partial sums. Hence the two actual value sets coincide
in this example. No value is assigned to any divergent series.

The formal term sequence is indexed from zero, with term number $j$
equal to $a_{j+1}2^{-(j+1)!}$ in the source indexing. Both this equality
and the surjectivity of the digit reindexing are proved, so no zeroth
factorial term has been inserted into the source series.

\section{Formal objects and verification}
\texttt{K} is Mathlib's \texttt{PadicAlgCl 2}, the algebraic closure with
the spectral norm extending the $2$-adic norm. \texttt{norm\_\allowbreak two}
uses that extension and the normalization on $\mathbf Q_2$.
\texttt{norm\_\allowbreak term\_\allowbreak of\_\allowbreak nonzero}
proves the lower bound (1), and the vanishing-term argument is in
\texttt{vanishing\_\allowbreak eventually\_\allowbreak zero}.
\texttt{series\_\allowbreak value\_\allowbreak rational} constructs the
finite rational sum explicitly and proves that its field embedding is
the actual sum value.

\texttt{sequential\_\allowbreak hasSum} and
\texttt{sourceSequentialValues\_\allowbreak eq} establish the ordinary
partial-sum bridge. The final theorem
\texttt{no\_\allowbreak continuum\_\allowbreak subfamily} directly
negates the existence of a subset of the source value set whose
\texttt{Cardinal.mk} equals $2^{\aleph_0}$. It therefore refutes a
necessary condition of the original claim without assuming anything
about the transcendence degree.

The project pins Lean 4.33.0, Mathlib and all dependency commits. A new
root project build, fourteen-declaration transitive axiom audit and
official fresh-kernel replay are recorded with source hashes. No
admitted proof, native decision shortcut or extra axiom is used.
The original bilingual source is preserved from upstream commit\par\noindent
\texttt{b355e42906e6eaf52bc7e9d005e0ef02d8aa7400}.\par
OpenAI Codex assisted the mathematics, proof and report. A separate
supervising Codex agent reviewed the sign, topology, indexing and both
convergence interpretations. This is AI review, not human certification.
The argument uses elementary necessary conditions for convergence and
countability; no new mathematical result is claimed.
\end{document}
